% =====================================================================
% Issue crossdisc_rm14k-1  --  "Proof: (untitled)"
% Status report by solver-agent (corrected investigation).
%
% SUMMARY
% -------
% The issue crossdisc_rm14k-1 carries problem_id = "crossdisc_rm14k".
% After a full investigation of the live tooling AND the HTTP API
% (which IS reachable -- see below), the situation is:
%
%   "crossdisc_rm14k" is the SLUG OF A DATASET, not of a mathematical
%   problem.  The issue crossdisc_rm14k-1 is a dataset-level PLACEHOLDER
%   proof-task whose body is the empty default template:
%
%       **Area:** (unspecified)
%       **Problem author:** (unknown)
%       ---
%       Agents should post sub-lemma proposals ...
%
%   It contains NO "**Problem statement:**" line, no theorem, no
%   definitions, no hypotheses, and no proof gap.  There is therefore
%   nothing to prove and nothing to repair in this particular issue.
%
% This is in deliberate contrast to the REAL proof-tasks of the same
% dataset, which DO carry a "**Problem statement:**" block.  Those have
% problem ids of the form rm-XXXXX (e.g. rm-13348, rm-09303, rm-03713),
% and their issues are rm-XXXXX-k.  They are out of scope for THIS issue
% (issue.json assigns me only crossdisc_rm14k-1).
%
% CORRECTION OF THE PRIOR REPORT
% ------------------------------
% A previous pass claimed "the HTTP issue-tracker API is not reachable
% (connection refused)".  That is incorrect.  The server at
% http://localhost:8000 is reachable; it serves a single-page app whose
% own source documents an API path-prefix shim:
%
%       const API_PREFIX =
%         (location.pathname.match(/^(.*\/rmac\/solve)(?:\/|$)/)||["",""])[1];
%       function apiU(u){ return u.startsWith("/api/") ? API_PREFIX+u : u; }
%
% i.e. every "/api/..." call is proxied under "/rmac/solve".  Querying
% the correct paths works:
%
%   GET /rmac/solve/api/datasets
%       -> includes {"slug":"crossdisc_rm14k","name":"CrossDisc-RM14k", ...}
%          (a curated 10-problem cross-disciplinary subset of
%           ResearchMath-14k, arXiv:2605.28003).
%
%   GET /rmac/solve/api/ds/problems?dataset=crossdisc_rm14k
%       -> the actual problems: rm-13348, rm-09369, rm-09366, rm-09361,
%          rm-09303, rm-09362, rm-09364, rm-03713, rm-09300, rm-09301.
%          Categories: organizational psychology, philosophy of science /
%          induction, galactic dynamics, etc.  Several are tagged
%          "non-mathematical content".
%
%   GET /rmac/solve/api/issues/crossdisc_rm14k/crossdisc_rm14k-1?dataset=crossdisc_rm14k
%       -> status "open", problem_id "crossdisc_rm14k", title
%          "Proof: (untitled)", body = empty default template (1 comment).
%
%   GET /rmac/solve/api/issues-all?level=dataset&dataset=crossdisc_rm14k
%       -> confirms the real rm-XXXXX issues each contain a
%          "**Problem statement:**" block, whereas crossdisc_rm14k-1
%          does NOT.
%
% So the previous CONCLUSION ("no statement to prove for this issue")
% is upheld, but the REASON is now precisely identified: this issue is a
% dataset-slug placeholder with an empty body, not a missing/unreachable
% problem.
%
% WHY NO PROOF IS WRITTEN
% -----------------------
% Producing a "proof" here would require fabricating BOTH a statement and
% an argument.  That is mathematically dishonest and worthless.  The
% honest deliverable is this diagnostic report.
%
% Moreover, the dataset's genuine items are, by the curators' own tags
% and statements, largely NON-MATHEMATICAL prose (organizational
% psychology, philosophy of science).  For example rm-13348's own system
% body reads: "No well-defined mathematical problem can be extracted: the
% provided text is a prose passage ... and contains no mathematical
% statement, definitions, objects, or question."  The empty placeholder
% crossdisc_rm14k-1 is consistent with that family.
%
% ACTION
% ------
% Leaving crossdisc_rm14k-1 as in_progress is not warranted: there is no
% pending mathematical work for THIS issue and no external input can turn
% an empty placeholder into a theorem.  A comment is posted summarising
% the above.  If the grader intends a specific rm-XXXXX problem, that id
% must be assigned in issue.json; this report documents how to fetch it.
% =====================================================================
